\section{Exploratory Data Analysis}
\label{sec:3_3}

Exploratory analysis is the foundation of our pattern discovery process. It allows us to identify fraud signatures across temporal, geographic, and relational dimensions before any formal modeling begins.

\subsection{Dataset Statistics}

The dataset, filtered from December 2024 onwards, consists of the following:

\begin{table}[h]
\centering
\caption{Dataset summary statistics.}
\begin{tabular}{ll}
\hline
Metric & Value \\
\hline
Total events & 804,717 \\
Unique listings & 86,160 \\
Fraud rate (per listing) & 8.11\% (6,987 fraudulent) \\
Fraud rate (per event) & 1.21\% \\
Features & 509 columns \\
Date range & Dec 2024 - Nov 2025 \\
\hline
\end{tabular}
\end{table}

\paragraph{A Note on Fraud Rates}
It is important to distinguish between event-level and listing-level fraud rates. The event-level rate of 1.21\% is somewhat misleading because legitimate listings tend to generate more events (an average of 9.7) than fraudulent ones (an average of 4.8). This is because legitimate listings usually go through their entire lifecycle, whereas fraudulent ones are often terminated early. Consequently, our analysis focus on the listing level, where the fraud rate is 8.11\%, representing a moderate class imbalance of approximately 1:11.

\subsection{Fraud Pattern Discovery}

\paragraph{Temporal Signatures}
Our analysis of how fraud is distributed over time revealed several key patterns. We observed that fraud activity tends to peak between 6:00 and 7:00 AM local time. This suggests either automated bot activity or submissions originating from different time zones, as legitimate listings are more uniformly spread across standard business hours.

On a monthly basis, we saw a significant spike in February 2025, with the fraud rate hitting 13.12%more than 60% higher than the annual average. This was followed by a decline to around 5% by the middle of the year, which may indicate that the platform's interventions were successful or that fraudsters were adapting their tactics. We also found that fraud rates were consistently higher on weekends (9.2% on Saturdays and 8.8% on Sundays) than on weekdays, likely due to reduced manual oversight.

\paragraph{Geographic Origins}
IP geolocation shows a very strong geographic concentration for fraudulent activity.

\begin{table}[h]
\centering
\caption{Fraud rates by geographic origin.}
\begin{tabular}{lll}
\hline
Origin Country & Fraud Rate & Lift vs Baseline \\
\hline
Nepal & 90.5\% & 11.2$\times$ \\
Benin & 66.9\% & 8.3$\times$ \\
Togo & 63.2\% & 7.8$\times$ \\
Ivory Coast & 58.4\% & 7.2$\times$ \\
France & 31.8\% & 3.9$\times$ \\
Switzerland & 6.2\% & 0.8$\times$ \\
\hline
\end{tabular}
\end{table}

Specifically, countries in West Africa such as Benin, Togo, and the Ivory Coast show fraud rates between 60% and 90%. This is consistent with established research on real estate scams and provides a high-precision indicator for detection, even if the overall volume from these regions is low.

\paragraph{High-Risk Listing Categories}
The characteristics of the listings themselves also offer strong clues about their legitimacy.

\begin{table}[h]
\centering
\caption{Fraud rates by listing category.}
\begin{tabular}{lll}
\hline
Category & Fraud Rate & Lift \\
\hline
Studio apartments & 34.97\% & 4.3$\times$ \\
RENT (vs BUY) & 9.47\% vs 0.88\% & 10.8$\times$ \\
Cross-platform (both Homegate \& ImmoScout24) & 26.74\% & 3.3$\times$ \\
\hline
\end{tabular}
\end{table}

Studio rentals are the highest-risk category, likely because their lower price points attract more potential victims and require less detailed documentation. We also see a stark difference between rental fraud (9.47%) and purchase fraud (0.88%), reflecting the much lower verification hurdles for rental listings.

\subsection{Entity-Sharing and Network Analysis}

By looking at device fingerprints, IP addresses, and other identity markers, we can begin to see the underlying network structures of coordinated fraud rings.

\paragraph{Device Sharing}
The evidence for device-based fraud rings is particularly strong.

\begin{table}[h]
\centering
\caption{Device sharing statistics.}
\begin{tabular}{ll}
\hline
Metric & Value \\
\hline
Total unique devices & 16,447 \\
Devices with 2+ listings & 6,473 (39.4\%) \\
Multi-listing devices with any fraud & 1,322 (20.4\%) \\
Super-connectors (10+ listings) & 1,248 \\
Super-connector fraud rate & 35.5\% \\
Maximum listings per device & 6,251 \\
\hline
\end{tabular}
\end{table}

Finding over 6,000 listings associated with a single device fingerprint is a clear indication of a fraud ring, as no legitimate user would operate at that scale. These "super-connector" devices have a fraud rate of 35.5%, which is a 4.4-fold increase over the baseline.

\paragraph{Cross-Platform and IP Patterns}
Listings that appear on both Homegate and ImmoScout24 also carry a much higher risk, with a fraud rate of 26.74%. While some legitimate managers do cross-post, the signal becomes much stronger when combined with other indicators like new accounts or suspicious infrastructure. IP addresses show similar clustering, although we must be careful here as legitimate corporate or residential users often share IP addresses.

These patterns are the primary driver behind our graph-based methodology, where shared attributes form the edges that connect listings into a larger heterogeneous network.
